\section{Conclusions and future work} \label{sec:conclusions}

Using Lean, it is possible to automatically verify the different 
lemmas that relate the semantics of the ShEx language variants. 
The resulting github repository~\footnote{\url{https://anonymous.4open.science/r/shex_lean-B1A5/}}
contains the ShEx schemas and the inheritance paper semantics and proofs,
all the examples from both papers associated with their corresponding numbers, 
an executable main with two command-line options: \lean{list} to show all 
 the propositions and examples that can be checked, 
 and \lean{run} to check a specific one by its name.
Some of the Lean proofs are complex and were defined with the help of 
 Claude Code showing that the use of generative AI tools can be useful 
 to help automate the proof construction process.

 We consider that this development can serve as a first approach to explore 
  more design decisions about the ShEx language. 
Future work includes a more extensive exploration of the 
 design space for RDF validation languages including extending ShEx 
 for RDF 1.2 or other language features~\cite{L24} as well as developing 
  a Semantic web library of lean tactics that can be used to shorten the proofs and 
  verify the correctness of other Semantic Web languages and algorithms. 



